\section{2006秋}
\subsection{微分積分}
\prob{
  微分方程式
  \begin{gather}
    \frac{d^2 y}{dx^2} + 3\frac{dy}{dx} - 4y =  -4x + 3
  \end{gather}
  において、境界条件
  \begin{gather}
    y(0) = 0, \quad \frac{dy}{dx}(0) = 1
  \end{gather}
  を満たす解を求めなさい。
}

\prob{
  \begin{enumerate}[label=(\arabic*), ref=(\arabic*), itemsep=5pt]
    \item $\disp y = e^{\sqrt{1 + x^2 + x^4} + 4}$に対して$\disp \frac{dy}{dx}$を求めなさい。
    \item $|x| = \log y$に対して$\disp \frac{dx}{dy}$を求めなさい。
  \end{enumerate}
}

\prob{
  次の不定積分を$t = \sqrt{1 + e^x}$とおくことにより求めなさい。
  \begin{gather}
    \int {e^{2x}}{\sqrt{1 + e^x}}dx
  \end{gather}
}

\prob{
  次の定積分を求めなさい。
  \begin{gather}
    \int_{0}^{1} \frac{x}{\sqrt{x + 1} - \sqrt{x}} dx
  \end{gather}
}


\newpage
\subsection{線形代数}
\prob{
  次の行列$\bA$の行列式$|\bA|$を求め、$|\bA| = 0$となる条件を求めなさい。
  \begin{gather}
    \bA = \begin{pmatrix}
      \gra & \beta & \grg & 1 \\
      \beta & \gra & \beta & 1 \\
      \grg & \beta & \gra & 1 \\
      1 & 1 & 1 & 0
    \end{pmatrix}
  \end{gather}
}


\prob{
  2次形式
  \begin{gather}
    ax^2 + bxy + cy^2
  \end{gather}
  に関する以下の問いに答えなさい。
  ただし、$a, b, c$はある定数である。
  \begin{enumerate}[label=(\arabic*), ref=(\arabic*), itemsep=0pt]
    \item この2次形式を$\bx^{\rm{T}}\bA\bx$と表すとき、2次の対称行列$\bA$を成分表示で与えなさい。
    ただし、$\bx = (x,y)^{\rm{T}}$であり、
    $(\cdot)^{\rm{T}}$は行列あるいはベクトル$(\cdot)$の転置を表すものとする。
    \item 座標変換
    \begin{gather}
      \begin{pmatrix}
        X \\[10pt] Y
      \end{pmatrix}
      =
      \begin{pmatrix}
       \disp \frac{1}{\sqrt{2}} & \disp -\frac{1}{\sqrt{2}} \\[10pt]
       \disp \frac{1}{\sqrt{2}} & \disp \frac{1}{\sqrt{2}}
      \end{pmatrix}
      \begin{pmatrix}
        x \\[10pt] y
      \end{pmatrix}
    \end{gather}
    を用いて、$a = 5,\,b = 2,\,c = 5$に対し、$\bx^{\rm{T}}\bA\bx$を$\bX = (X,Y)^{\rm{T}}$
    に関する2次形式に書き直し、この2次形式が正値であるかどうかを調べよ。
  \end{enumerate}
}

\prob{
  次の連立方程式を解きなさい。ただし、$C$はある定数である。
  \begin{gather}
    \begin{pmatrix}
      1 & -2 & 3 \\
      3 & -5 & 5 \\
      1 & -1 & -1
    \end{pmatrix}
    \begin{pmatrix}
      x \\ y \\ z
    \end{pmatrix}
    =
    \begin{pmatrix}
      6 \\ C \\ -4
    \end{pmatrix}
  \end{gather}
}
